\begin{solapartat}
\[ B(x) = \frac{\mu_0 I}{2}\,\frac{R^2}{(x^2+R^2)^{3/2}} \]

\punts{0,6} $B(x = 0) = \dfrac{\mu_0 I}{2R}$

\punts{0,4} $B(x = 0) = \dfrac{4\pi 10^{-7}\times 10}{2\times 0,05} = 1,26\times 10^{-4}\ \mathrm{T} = 126\ \mathrm{\mu T}$
\end{solapartat}

\begin{solapartat}
\punts{1} Si el camp magnètic al centre de l'espira esta dirigit cap al sentit positiu de l'eix $x$, $(B\vec{\imath}\,)$, el sentit del corrent elèctric és antihorari (justificació amb la regla de la mà dreta).

\figura[0.3]{sol_fig1.png}
\end{solapartat}
